% =============================================================================
% SECTION 3 — METHOD (revised for paper_v2)
% Forward references:
%   Section~\ref{sec:exp:ablation}, Section~\ref{sec:method:cot}
%   Appendix~\ref{app:data}, Appendix~\ref{app:algo}
% =============================================================================

\section{Method}
\label{sec:method}
\vspace{-0.5ex}

\subsection{Motivation: Function Calls as Agent-Like Structures}
\label{sec:method:motivation}

At step $t$ a coding agent maintains a history $h_t$, samples an action
$a_t \sim \pi(a_t \mid h_t)$, receives an observation
$o_{t+1} \sim p(o_{t+1} \mid h_t, a_t)$, and continues conditioned on the
trace. Under the function-call/agent-step correspondence
(Figure~\ref{fig:teaser}, left), the four parts---history, action, return,
continuation---align respectively with pre-call context, the call,
the returned value, and downstream usage. This isomorphism motivates
extracting an agentic conditioning signal directly from source code via a
fill-in-the-middle (FIM) objective. The correspondence is
\emph{structural} rather than literal---FIM conditions on a given suffix
while agent inference generates one---but a model that reasons
bidirectionally about function-level dependencies should learn a
representation that transfers when the suffix is generated rather than
observed. Standard random-span FIM~\citep{qwencoder, deepseekcoder,
starcoder} captures this structure only incidentally; we construct a
\emph{function-aware} variant whose masking targets are chosen by program
structure and contextual predictability.

\vspace{-0.5ex}
\subsection{Data Collection and Decontamination}
\label{sec:method:data}

We curate a corpus of $968$ Python repositories from GitHub. Starting
from $\sim\!2{,}000$ candidates retrieved by combining a star-count
threshold with ten topic queries, we apply manual quality filtering and
remove every repository whose origin overlaps with the source
repositories of SWE-Bench~\citep{swebench} (verified by repository name
and any known fork). To eliminate test-time leakage, we restrict each
repository to commits whose timestamp precedes the earliest base-commit
used in SWE-Bench-Verified and SWE-Bench-Lite. The filtered corpus
contains $\approx\!78\mathrm{K}$ self-contained Python files (averaging
$428$ lines per file) and yields $\sim\!2.6\mathrm{B}$ training tokens
under the Qwen2.5-Coder tokenizer, spanning $320\mathrm{K}$
single-function, $60\mathrm{K}$ pair, and $20\mathrm{K}$ triple FIM
samples; full statistics, topic-category breakdown, and license
inventory are reported in Appendix~\ref{app:data}, and the full
repository list will be released to support reproduction.

\vspace{-0.5ex}
\subsection{Function-Aware FIM Target Selection}
\label{sec:method:selection}

Given each source file, our pipeline produces a set of mask targets
together with the corresponding masked file. Each target is a single
function or a connected group of $2$--$3$ functions; the multi-function
variant is studied in Section~\ref{sec:exp:ablation}. The pipeline has
four stages: dependency-graph construction, complexity scoring,
inferability scoring, and threshold-based selection. Figure~\ref{fig:pdg_score}
illustrates these stages on a small running example---a
\texttt{Calculator} class with two top-level helpers. The full
single-function selection algorithm is given in
Appendix~\ref{app:algo:single}, the multi-function extension in
Appendix~\ref{app:algo:groups}, and a numerical walkthrough of every
quantity shown in Figure~\ref{fig:pdg_score}(b) in
Appendix~\ref{app:algo:worked_example}.

% -----------------------------------------------------------------------------
% Figure: PDG construction and H/I score schematic (compact single-row layout)
% -----------------------------------------------------------------------------
\begin{figure}[t]
\centering
\includegraphics[width=\linewidth]{figs/pdg_score.pdf}
\vspace{-2ex}
\caption{Function-aware FIM target selection on a small calculator
example. \textbf{(a)}~Program dependency graph parsed from the AST:
solid arrows are call edges $\mathcal{E}_{\mathrm{call}}$, dashed lines
are sibling edges $\mathcal{E}_{\mathrm{sib}}$ between same-class
methods. \textbf{(b)}~Stacked bars decompose the complexity score
$\hat{H}\!=\!0.40$ (Eq.~\ref{eq:complexity}; LoC, CC, depth) and the
inferability score $\hat{I}\!=\!0.48$ (Eq.~\ref{eq:inferability}; five
context signals) for $\mathtt{Calculator.total}$, yielding
$\mathrm{FIM}\!\approx\!0.22\!\geq\!\tau$. \textbf{(c)}~In the
$(\hat{H},\hat{I})$ plane, the four short helpers fall outside the
selected region while $\mathtt{Calculator.total}$ is kept.}
\label{fig:pdg_score}
\vspace{-1.5ex}
\end{figure}

\subsubsection{Program Dependency Graph}
\label{sec:method:pdg}

For each file we parse its AST and extract the set $\mathcal{V}$ of
function nodes (top-level functions and class methods, identified by
qualified names). Two edge sets are constructed: \emph{call edges}
$\mathcal{E}_{\mathrm{call}}$ between caller and callee, and
\emph{sibling edges} $\mathcal{E}_{\mathrm{sib}}$ between methods of the
same class (capturing intra-class coupling that does not manifest as
direct calls but carries information through shared instance state).
Call resolution handles common Python idioms (direct invocations, class
instantiation, \texttt{self}/\texttt{cls} method calls) with a
short-name fallback against the qualified-name index;
Appendix~\ref{app:algo:pdg} gives the full procedure.

\subsubsection{Complexity Score $\hat{H}$}
\label{sec:method:complexity}

For each $v \in \mathcal{V}$ we define
\vspace{-0.5ex}
\begin{equation}
  \hat{H}(v) \;=\; w_{\ell}\,\phi\!\left(\mathrm{LoC}(v),\,c_{\ell}\right)
            + w_{c}\,\phi\!\left(\mathrm{CC}(v),\,c_{c}\right)
            + w_{d}\,\phi\!\left(\mathrm{D}(v),\,c_{d}\right),
  \label{eq:complexity}
\end{equation}
where $\mathrm{LoC}(v)$ is lines of code, $\mathrm{CC}(v)$ is McCabe
cyclomatic complexity, $\mathrm{D}(v)$ is the maximum nesting depth of
control-flow constructs, and $\phi(x,c)\!=\!\min(x/c,\,2)$ normalizes
each quantity by a soft cap. Caps and weights are listed in
Appendix~\ref{app:algo:hyperparams}.

\subsubsection{Inferability Score $\hat{I}$}
\label{sec:method:inferability}

A target should be \emph{recoverable} from the surrounding context.
$\hat{I}(v)$ aggregates five context-derived signals that approximate
the mutual information between $v$'s body and the rest of the file:
\vspace{-0.5ex}
\begin{equation}
  \hat{I}(v) \;=\; \alpha\,C_{\mathrm{caller}}(v)
              + \beta\,C_{\mathrm{callee}}(v)
              + \gamma\,C_{\mathrm{sig}}(v)
              + \delta\,C_{\mathrm{doc}}(v)
              + \varepsilon\,C_{\mathrm{class}}(v).
  \label{eq:inferability}
\end{equation}
$C_{\mathrm{caller}}$ scores call-site argument specificity,
$C_{\mathrm{callee}}$ counts intra-file functions called by $v$,
$C_{\mathrm{sig}}$ aggregates type annotations and name descriptiveness,
$C_{\mathrm{doc}}$ indicates docstring presence, and
$C_{\mathrm{class}}$ counts in-class siblings; full component formulas
and weights are in Appendix~\ref{app:algo:inferability}. Each component
is a hand-designed proxy: a learned predictability score would couple
selection to a particular reference model and complicate the
cross-base-model generalization analysis (Section~\ref{sec:exp:main}).

\subsubsection{Single-Function Score}
\label{sec:method:single}

We combine $\hat{H}$ and $\hat{I}$ in a harmonic-mean-like form, scaled
by a one-sided difficulty penalty $\rho(\Delta(v))$:
\vspace{-0.5ex}
\begin{equation}
  \mathrm{FIM}(v) \;=\;
  \frac{\hat{H}(v)\,\hat{I}(v)}{\hat{H}(v) + \hat{I}(v) + \epsilon}
  \cdot \rho\!\left(\Delta(v)\right).
  \label{eq:fim_single}
\end{equation}
The harmonic-mean form forces both $\hat{H}$ and $\hat{I}$ to be large
simultaneously, penalizing imbalance; $\rho$ down-weights targets that
remain hard even given full context, which would otherwise be
unlearnable noise. Hard filters on length and dunder methods, the form
of $\rho$, and the threshold $\tau\!=\!0.20$ used throughout this work
are documented in Appendix~\ref{app:algo:penalty}.

\subsubsection{Multi-Function Group Selection}
\label{sec:method:groups}

Real-world code patches frequently span multiple related
functions~\citep{swebench}, motivating an extension that masks
\emph{groups} of $k\!=\!2$ or $k\!=\!3$ structurally connected functions;
we use this variant in our main recipe and ablate it in
Section~\ref{sec:exp:ablation}. A group score $\mathrm{FIM}(G)$
multiplies a coupling term, the harmonic-mean-like product of
group-level $\hat{H}(G)$ and $\hat{I}(G)$, and a difficulty penalty,
with $\hat{I}(G)$ recomputed under joint masking so that intra-group
references cannot inflate the score. The eight topology patterns
(caller-callee, co-callee, sibling-coupled, mutual-call, call-chain,
hub, fan-in, class-triad), the full equations,
Algorithm~\ref{alg:selection_multi}, and a worked pair example are in
Appendix~\ref{app:algo:groups}.

\vspace{-0.5ex}
\subsection{Chain-of-Thought Augmentation}
\label{sec:method:cot}

For each selected target we generate a natural-language rationale via
Gemini-3-Preview prompted with the masked file and the original function
body (template in Appendix~\ref{app:algo:cot}). The rationale is
embedded \emph{inside} the FIM middle span, immediately before the
function body:
\vspace{-0.5ex}
\begin{center}\small
\texttt{<fim\_prefix>}\,$\langle$prefix$\rangle$\,%
\texttt{<fim\_suffix>}\,$\langle$suffix$\rangle$\,%
\texttt{<fim\_middle>}\,$\langle$rationale$\rangle$\,$\langle$body$\rangle$
\end{center}
The model is therefore trained to produce, in order, a reasoning trace
followed by code consistent with that trace, mirroring the
think-then-act structure of an agent step. To control quality, each
$(\text{CoT}, \text{completion})$ pair is then scored by Gemini-3-Flash
given the masked file, the rationale, the generated completion, and the
ground-truth function body; we retain only the top-scoring
$\sim\!400$K samples as our final training corpus. The verification
prompt template is in Appendix~\ref{app:algo:cot}.
Section~\ref{sec:exp:ablation}
isolates the contribution of FIM structure from CoT distillation via a
\emph{self-CoT} variant in which the model under training generates its
own rationales.
